\documentclass[11pt]{article}
\usepackage{amsmath,amssymb,amsthm}
\usepackage[margin=1in]{geometry}
\newtheorem{lemma}{Lemma}
\title{Room of Mathematicians, item 5 (P5, $W_T$ non-$D$-finite) --\\
Seat: Hilbert -- search for any meromorphic continuation of $F_T$}
\author{}
\date{}
\begin{document}
\maketitle

\section*{Calibration point}

Paper1 (\S8.8 of \texttt{paper/journal/paper1.tex}, around line 2317
for the Proposition and line 2548 onward for $F_T$) sets up $W_T$ as follows.
$F_T(x)=\sum_{d\ge0}u_d^Tx^d$ is the spherical growth series of $W_T\le
\mathrm{Isom}(\mathbb R^2)$, the side-reflection group of a rational right
triangle $T$, with respect to the three reflection generators
$\{R_0,R_1,R_2\}$: $u_d^T$ counts elements of word length exactly $d$.
Every element of $W_T$ is represented by an \emph{edge profile}
$A=\sum_j a_jt^j\in\mathbb Z[t^{\pm1}]$ (a formal Laurent polynomial recording,
at each site $j$ of the one-dimensional ``lamp'' walk that a geodesic word
induces, the signed deposit made there), and two profiles $A,A'$ represent the
\emph{same} element of $W_T$ iff $A-A'=2\mu_Tf$ for some $f\in\mathbb
Z[t^{\pm1}]$, where $\mu_T$ is $\zeta_T$'s minimal polynomial and $f$ is called
a \emph{carry}. Evaluating a profile at $t=\zeta_T$ (a point on the unit
circle, root of the primitive self-reciprocal quadratic $ct^2-et+c$) gives the
actual translation vector in $\Lambda_T=2(\zeta_T-1)A_T\subset\mathbb C$, and
$\zeta_T$ itself is the linear part of the ``rotation'' generator $r=R_xR_h$.

So $F_T$ counts orbits of profiles under $A\mapsto A+2\mu_Tf$, weighted by the
geodesic length of the group element each orbit represents. Concretely: to
express a given target profile $A$ using the smallest possible word, one
performs a division-with-remainder process, subtracting multiples of
$2\mu_T$ from $A$'s low-degree end, exactly as in a positional numeral system
with ``base'' $\zeta_T$ -- except that $|\zeta_T|=1$, so this base does not
contract. In an ordinary numeration algorithm with an expanding/contracting
base $b$, $|b|>1$, the running remainder (the ``carry'') stays in a bounded
window because each division step shrinks it geometrically; here nothing
shrinks it, and the carry state, tracked as a point of the lattice
$\mathbb Z[\zeta_T]\subset\mathbb C$ (i.e.\ literally a point of the plane,
not a bounded digit), performs an unbounded walk as $d$ grows. This is the
``planar carry'': the state of the numeration/carry process used to compute
$W_T$-geodesic length is not a bounded automaton state but a genuinely
two-dimensional, apparently unbounded, walk. Paper1's own numerics
(\texttt{wt\_growth/wt\_carry}, shape $(c,e)=(2,1)$) show the carry width
needed to realize geodesics of the universal ball of radius $d$ growing like
$\lfloor(d-12)/2\rfloor$, i.e.\ linearly, for all sampled $d\le32$ -- verified
in that range, not proved in general.

This reproduces, in my own words, exactly the mechanism paper1 and
\texttt{P5\_note.tex}/\texttt{P5\_v2\_note.tex} describe; no correction to it
is being made here.

\section*{What was tried (different from the kernel-method route already
attempted in P5\_v2)}

P5\_v2 already tried, and found blocked, the kernel-method / catalytic-variable
approach: it needs a functional equation of the form ``sum over the last step
of a Laurent polynomial in a catalytic variable $v$ tracking the running
carry value,'' and found that an $a^{\pm1}$-step at lamp-height $h$ forces $v$
to be raised to the \emph{non-integer, $h$-dependent} exponent $\zeta_T^{\,h}$
-- not a valid catalytic step for a formal-power-series kernel method (the
method requires the exponent shift to be a fixed integer or fixed Laurent
polynomial, independent of the step index). I did not re-attempt this; I
looked for a genuinely different route into continuation of $F_T$, along the
four lines the task specifies.

\paragraph{(1) Integral/contour representation.} The natural candidate would
be a Cauchy-type integral extracting, for each $d$, the number of
carry-inequivalence classes of length exactly $d$, e.g.\ something of the
shape
\[
u_d^T \;=\; \frac{1}{2\pi i}\oint \Phi(z,\zeta_T)\,\frac{dz}{z^{d+1}}
\]
for some kernel $\Phi$ built from the carry lattice. This requires knowing, in
closed form, the indicator/weight of ``profile $A$ is $2\mu_T$-reducible to a
given minimal-length representative'' as a function that can be Fourier/Cauchy
expanded. For the travel operator behind $U,V,G^T_0$ this indicator comes from
an explicit finite-state transfer matrix (the travel block), which is why a
Fredholm-determinant representation exists there. For $F_T$ the corresponding
``reducibility'' relation is division by $2\mu_T$ in $\mathbb Z[t^{\pm1}]$
evaluated on the unit circle: this is long division in a Laurent-polynomial
ring, not evaluation of a finite automaton, and I could not construct any
kernel $\Phi$ whose contour integral reproduces even the first few $u_d^T$
without already knowing the full carry-equivalence classes -- i.e.\ the
integral representation I could write down is a tautological restatement of
the counting problem (integrate an indicator function that is exactly as hard
to describe as $F_T$ itself), not a genuine transform of a simpler object.
This is not a proof that no such representation exists; it is a report that
none was found, and a diagnosis of why the obvious attempt is circular: unlike
a finite-state transfer, division by $2\mu_T$ does not admit a compact
generating-function kernel because the remainder ranges over an infinite,
non-compact index set (Gap 2 already isolates exactly this).

\paragraph{(2) Special-function connection (theta / modular / $q$-Bessel).}
The $q$-Bessel functions used elsewhere in this program (for the (R-J) /
travel-pole analysis of $U,V$) arise from a genuinely different structure: a
recursion in a variable that is being sent to $1$ along a specific limiting
regime (the confluent/travel limit), with the q-difference structure coming
from the finite-state travel block. $F_T$'s carry process has no such
recursion in a small parameter -- $\zeta_T$ is a \emph{fixed} unit-circle
algebraic number, not a deformation parameter, and the ``base'' does not
degenerate to a nice limit. Theta functions and canonical modular objects
typically arise from lattice sums $\sum_{\lambda\in L}q^{Q(\lambda)}$ for a
contracting or at least well-behaved quadratic form $Q$ and a genuine radius
of convergence in $q$; here the natural lattice sum would be over carry values
$f\in\mathbb Z[\zeta_T^{\pm1}]$ weighted by geodesic length, but the weight
(geodesic length as a function of the carry) is not given by a quadratic form
or any other closed algebraic expression -- it is itself defined by an
optimization (the min-plus DP that \texttt{wt\_carry} runs), which is exactly
the ``no functional equation is known'' obstruction (Gap 3), not something a
change of special-function variable bypasses. I did not find a substitution
turning $F_T$ into a theta or $q$-Bessel object, and I do not believe one is
likely to exist without first solving the min-plus optimization in closed
form, which is the same open problem in different language.

The one genuine structural fact available (Proposition~2, part (v) of
paper1, ``Structure of $W_T$'') is that $H_T\le W_T$ is a cocompact lattice in
$G_T=(\mathbb C\times\prod_{\pi\mid c_T}\mathbb Q(i)_\pi)\rtimes_{\zeta_T}
\mathbb Z$, so $W_T$ is quasi-isometric to $\mathbb R^2$ times a horocyclic
product of Bruhat--Tits trees (for $T=(1,2)$: the Diestel--Leader graph
$\mathrm{DL}(5,5)$). This is worth recording as a considered-and-rejected
route: growth series of $\mathrm{DL}(q,q)$ \emph{as a vertex-transitive graph}
(equivalently, of the lamplighter group $\mathbb Z_q\wr\mathbb Z$ with its
standard generators) are known to be rational (Bartholdi--Neuhauser--Woess-type
results on lamplighter growth). But $H_T$'s lamp group here is $A_T\cong
\mathbb Z[\zeta_T^{\pm1}]$, of rank $2$ over $\mathbb Z[1/c_T]$ and \emph{not}
finite and \emph{not} finitely generated (paper1 Prop.\ (i)); the
quasi-isometry to $\mathrm{DL}(5,5)$ is a coarse statement (Milnor--Švarc) that
does not preserve the exact word-length metric or its generating function. So
the known rational growth series of the finite-lamp Diestel--Leader graphs
does not transfer, and paper1's own Remark on scope already flags that $W_T$
is emphatically not an instance of the finite-lamp $\mathbb Z\wr\mathbb Z$
picture. I record this as a route checked and rejected, not a new obstruction.

\paragraph{(3) Coefficients against known sequences.} By universality
(Theorem \texttt{thm:universality-sharp}), $u_d^T=u_d$ (the shape-independent
universal sequence, OEIS A396406) for all $d\le n_T$, and only diverges from it
beyond the shape-dependent onset $n_T\ge c_T$. So the early coefficients of
$F_T$ for any specific $T$ are, by construction, not a new sequence to check
against OEIS -- they coincide with the already-catalogued A396406 until the
onset, and the repository's own numerics (\S ``Numerical record'' in
paper1.tex, and \texttt{wt\_growth/}) already track the post-onset deficits
$u_d-u_d^T$ for $(1,2)$ beyond depth 33. I read that table; it shows no
periodicity or evident closed form in the deficits at the depths recorded, and
Theorem~\texttt{thm:WT-lowcomplexity} already rules out rational/algebraic/
P-recursive relations up to size boxes computed there for \emph{all} $T$ at
once (using exactly the shape-independent layers). I did not re-run
\texttt{wt\_growth} or \texttt{guess.py} for this note (both are read-only per
task instructions and re-running would not change coefficients already
certified); no hidden closed form suggests itself from re-reading this
existing data.

\paragraph{(4) Why continuation is hard here specifically.} The precise
obstruction, sharper than ``nobody has tried'': continuation via pole
confinement (the route that works for $U,V,G^T_0$) needs an explicit operator
or transfer equation whose resolvent/Fredholm determinant \emph{is} the
generating function, so that poles are visible as zeros of that determinant
inside its domain of analyticity. That representation exists for $U,V,G^T_0$
because their combinatorics come from a \emph{finite-state} travel block: a
bounded catalytic variable, damped or at least confined, admitting a
finite-dimensional transfer matrix whose characteristic/Fredholm data give
meromorphic continuation for free. For $F_T$, the analogous catalytic
structure -- the carry $f$ -- is proved (paper1's own numerics, and the
$h$-dependent exponent $\zeta_T^{\,h}$ obstruction found independently by
P5\_v2) to be an \emph{unbounded, genuinely two-dimensional} state, and worse,
it enters the length function non-linearly (as the outcome of a min-plus
optimization, not a fixed weight per carry value). Both of these facts block
the standard routes to continuation, for related but distinct reasons:
\begin{itemize}
\item An unbounded catalytic variable defeats the finite-state/kernel-method
machinery (this is Gap 2/Gap 3, and is what P5\_v2's non-integer-exponent
computation makes precise).
\item A non-fixed (optimization-defined) weight defeats any na\"ive lattice-sum
or theta-function ansatz, because such ans\"atze need the weight of a lattice
point to be given in closed form \emph{before} summing, whereas here the
weight (geodesic length) is only known as the value of a dynamic program.
\item The base being exactly on the unit circle ($|\zeta_T|=1$) is also the
reason no results from contracting-base numeration systems (canonical number
systems, Akiyama--Frougny--Sakarovitch rational-base systems, all of which
require $|{\rm base}|>1$) apply even qualitatively: those theories all use
contraction to bound the carry, and $|\zeta_T|=1$ is exactly the boundary case
they exclude, not a variant of a case they cover.
\end{itemize}
None of these three points is new mathematics: they restate, in the
integral/special-function language the task asked me to try, exactly the same
Gap 2/Gap 3 obstruction paper1 and \texttt{P5\_note.tex}/\texttt{P5\_v2\_note.tex}
already isolate. What is added here is the explicit report that the two most
natural \emph{different-in-form} attempts (a Cauchy/contour kernel, and a
theta/$q$-Bessel substitution) both run into the identical obstruction from a
new angle, rather than a new gap.

\section*{Verdict}

\textbf{NOTHING WORKED.} No integral representation, no special-function
identification, and no hidden closed form in the computed coefficients was
found. Precisely why: any route to a meromorphic continuation of $F_T$ that
mimics the $U/V/G_0^T$ route needs a finite-dimensional (or otherwise tame)
transfer structure carrying the analytic information; the actual combinatorial
engine of $F_T$ is division by $2\mu_T$ in $\mathbb Z[t^{\pm1}]$ evaluated at a
unit-circle point $\zeta_T$, whose remainder (``carry'') is an unbounded planar
walk entering the length function only through a min-plus optimization, not a
fixed per-symbol weight. This is a structural mismatch with every method tried
(kernel method in P5\_v2; contour/Cauchy kernel and theta/$q$-Bessel
substitution here), not merely an unexplored direction, and it matches exactly
what Gaps 1--3 of \texttt{P5\_note.tex} already say is missing. This note adds
no new proved lemma and does not claim otherwise; it records two additional
failed attempts and a sharper (still informal) diagnosis of the common
obstruction, consistent with -- and not contradicting -- the existing gap
statements.

\end{document}
